% 01_body.tex — 산출물 1/5 비즈니스 케이스 요약 (빈 템플릿 / 완성 예시 공용 본문)
% 드라이버: 01_비즈니스케이스요약_빈템플릿.tex (\wbblanktrue) / 01_비즈니스케이스요약_완성예시.tex (\wbblankfalse)
% 내용 출처: PM트랙/Day1_워크북.md §1.3(항목 가이드) §1.4(빈 템플릿) §1.5(온담 P1 완성 예시본)
\documentclass[10pt,a4paper]{article}
\usepackage[a4paper,margin=16mm,top=14mm,bottom=20mm]{geometry}
\def\DocNum{1}
\def\DocTitle{비즈니스 케이스 요약}
\def\DocSub{Business Case Summary}
\def\DocVariant{\B{빈 템플릿}{온담 P1 완성 예시}}
\usepackage{wbstyle}
\begin{document}

\docheader
 {\Bg{[정식 명칭과 코드]}{차세대 주문·재고 통합 플랫폼 구축 (ONE ONDAM P1)}}
 {\Bg{[v0.x / 기준 시점]}{v0.9 (G0 착수용) / 2026-01-05}}
 {\Bg{[PM 실명]}{P1 PM (발주사 IT본부)\rd{7mm}}}
 {\Bg{[스폰서 또는 재무 담당 실명]}{정미란 CFO (프로그램 스폰서)\rd{7mm}}}

% ------------------------------------------------------------------ 1
\secbar{1. 문서 정보}
\guide{이 요약본이 어느 원본 승인 문서를 근거로 하는지, 몇 번째 버전인지, 어느 시점의 숫자인지를 적는다. 게이트마다 갱신되므로 버전·기준 시점이 없으면 어느 숫자가 현행인지 알 수 없다.}
{\footnotesize
\begin{tabularx}{\linewidth}{|L{32mm}|Y|}\hline
\hd{항목} & \hd{내용} \\ \hline
\B{\rd{8mm}}{}프로젝트명 & \Bg{[정식 명칭과 코드]}{차세대 주문·재고 통합 플랫폼 구축 (ONE ONDAM P1)} \\ \hline
\B{\rd{8mm}}{}원본 승인 문서 & \Bg{[부의안·품의서 명칭, 결재번호, 승인일]}{2025-12 이사회 의결 「ONE ONDAM 프로그램 추진(안)」 및 별지 편익 산정표 (이사회 부의안 2025-12)} \\ \hline
\B{\rd{8mm}}{}버전 / 기준 시점 & \Bg{[v0.x / 숫자를 뽑은 회계기간]}{v0.9 (G0 착수용) / 2025년 연결 재무제표, 2025년 운영 KPI} \\ \hline
\B{\rd{8mm}}{}작성자 / 검토자 & \Bg{[PM 실명 / 스폰서 또는 재무 담당 실명]}{P1 PM (발주사 IT본부) / 정미란 CFO(프로그램 스폰서), 재경팀장} \\ \hline
\B{\rd{8mm}}{}다음 갱신 시점 & \Bg{[게이트명과 날짜]}{G1 기획 완료 2026-04-30 — 원가·일정 기준선 확정과 함께 v1.0} \\ \hline
\end{tabularx}}

% ------------------------------------------------------------------ 2
\secbar{2. 전략적 배경 (Strategic Case) — 왜 지금인가}
\guide{조직이 처한 상황과 해결하려는 문제. 3~5문단, 문단마다 숫자 1개 이상, 측정값의 분모를 밝힌다. 마지막 문단은 반드시 ‘하지 않을 경우(do-nothing)’의 시나리오. 솔루션을 배경으로 적지 말 것.}
\begin{fieldbox}
\ifwbblank
\lines{8}{9.5mm}
\else
온담F\&B홀딩스는 2025년 연결매출 8,740억 원(C-01)으로 3년간 18\% 성장했으나 영업이익은 371억에서 312억으로 16\% 감소했고, 영업이익률은 5.0\%에서 3.6\%로 내려앉았다(C-02). 부채비율 187\%(C-04), 이자보상배율 2.69배(C-06)인 회사가 428억 원의 프로그램을 승인한 것은 여유가 있어서가 아니라 세 가지 압력이 동시에 도달했기 때문이다.

첫째, 기술 부채가 한계에 이르렀다. 615개 매장(C-09)의 결제를 처리하는 OD-POS 2.1은 2011년 자체개발 시스템으로 서버 OS 제조사 지원이 2024년 7월에 종료되었고 연장지원 계약도 없다. 원 개발자 3명 중 재직자는 1명이다. SAP ECC 6.0의 표준 유지보수는 2027년 12월에 끝난다. 시스템 6종을 잇는 인터페이스 47본 중 13본은 설계문서가 없다(C-18). 어느 쪽이든 2027년 안에 손을 대야 한다.

둘째, 운영 손실이 측정 가능한 규모로 누적되고 있다. 폐기율 3.9\%는 식자재 취급원가 2,530억(C-15) 기준 연 약 98.7억 원의 식자재가 버려진다는 뜻이다. 결품률 4.7\%, 재고 정확도 91.3\%(순환실사 SKU 69\%만 분모, O-08), 가맹 발주 리드타임 38시간(O-12), 가맹 발주의 39\%가 전화·팩스·엑셀 메일로 들어와 본사 담당자가 SAP에 수기 입력한다. 온라인 매출 비중은 6.2\%(O-03)로 경쟁 브랜드 대비 낮고, 배달3사 경유 매출 242.7억에 수수료율 12.8\%(O-04, O-05)가 붙는다.

셋째, 외부 시계가 고정되어 있다. 2019년 350억을 투자한 사모펀드의 2027년 IPO 또는 매각 조항에 따라 회사는 2027년 3분기 IPO 예비심사 청구를 목표로 하며, 개인정보 보호법 개정 대응(2026-09-11), 전자금융업 등록 심사(2027-03-31), ISMS-P 인증 심사(2027-06)는 협상할 수 없는 날짜다.

\textbf{하지 않을 경우(do-nothing).} 보안 패치가 나오지 않는 OS 위에서 615개 매장의 결제를 계속 처리한다. 매장 결제 시스템의 보안 사고는 IPO 예비심사와 ISMS-P 심사에 직접적 결격 사유가 되며, 개인정보 처리항목 142개·위탁사 19개·재위탁 4건을 현행 아키텍처에서 규제 기한 안에 정비하는 것은 IT본부 47명(C-11)의 역량으로 불가능하다는 것이 IT본부의 판단이다. 폐기 손실 98.7억과 수기 입력 인건비는 그대로 남는다. 즉 do-nothing은 “비용 0”이 아니라 “규제 결격 + 연 100억 규모 손실 지속”이다.
\fi
\end{fieldbox}

% ------------------------------------------------------------------ 3
\secbar{3. 대안 검토 (Options Appraisal)}
\guide{최소 세 대안 — 아무것도 안 함 / 최소한만 함 / 제안안 — 과 기각된 대안 1개 이상. 비교 기준(예: 3년 TCO)을 통일하고, 기각 사유는 한 줄이라도 반드시 적는다.}
{\footnotesize
\begin{tabularx}{\linewidth}{|L{22mm}|Y|L{20mm}|Y|Y|Y|}\hline
\hd{대안} & \hd{개요} & \hd{개략 비용(억)} & \hd{주요 편익} & \hd{주요 리스크} & \hd{채택 여부 / 사유} \\ \hline
\ifwbblank
\rd{22mm}0. 아무것도 안 함 & \g{[현행 유지 시 무엇이 지속되는가]} & \g{[기회비용 포함]} & & & \g{[기각 — 사유]} \\ \hline
\rd{22mm}1. 최소한만 함 & & & & & \\ \hline
\rd{22mm}2. \g{[기각된 대안]} & & & & & \\ \hline
\rd{22mm}3. 제안안 & & & & & \g{[채택 — 사유]} \\ \hline
\else
0. 아무것도 안 함 & 현행 6개 시스템 유지, OS 연장지원 미계약 상태 지속 & 0 (기회비용 연 ~100) & 없음 & 보안 사고·규제 결격·IPO 결격 & 기각 — 규제 기한 충족 불가 \\ \hline
1. 최소한만 함 & OD-POS 서버 OS 마이그레이션 + 인터페이스 13본 문서화 + 개인정보 처리항목 정비만 수행 & 약 18 [추가 설정] & 규제 기한 충족, 보안 노출 해소 & 실시간 재고·옴니채널 편익 0, 2027-12 SAP 종료 시 재투자, 원 개발자 1명 의존 지속 & 기각 — 편익 없이 2027년 재투자 불가피. 단, 제안안 지연 시 폴백(fallback)으로 유지 \\ \hline
2. 상용 리테일 패키지 도입 & 글로벌 리테일 POS·OMS 패키지를 도입하고 SAP과 표준 연동 & 약 140~190 (라이선스 비중 높음) [추가 설정] & 구축 기간 단축 가능성 & 직영·가맹·B2B 세 사업 모델을 한 패키지가 다루지 못함, 가맹 발주·B2B EDI 커스터마이징 대량 발생, 벤더 종속(IT본부 운영 불가) & 기각 — 2022년 WMS 사례처럼 벤더 소유 데이터 구조로 인한 범위 축소 재발 우려 \\ \hline
3. 통합 플랫폼 SI 구축 (제안안) & 클라우드 POS + 주문 허브 + 통합 재고 + 가맹 포털 + 통합 계층·데이터 전환. SAP ECC는 유지·연동 & 168.0 (P-08) & §4 참조 & 17개월 일정, P2 의존, 3PL 연동 협상, 가맹 388점 동의 & \textbf{채택} — 세 사업 모델을 하나의 주문·재고 원장으로 다루는 유일한 대안이며 IT본부 자체 운영 요건 충족 \\ \hline
\fi
\end{tabularx}}

% ------------------------------------------------------------------ 4 (가로)
\landscapeon
\secbar{4. 편익 (Benefits) — 무엇이 얼마나 좋아지는가}
\guide{편익 항목별 현재 기준값·목표값·달성 시점·연간 금액·산출 분모·측정 책임자. \textbf{분모 열과 책임자 열이 없는 편익 표는 편익 표가 아니다.} 이 프로젝트가 직접 기여하는 항목과 타 프로젝트 의존 항목을 구분하고, 금액 환산이 안 되면 ‘미산정(사유)’으로 정직하게 쓴다. 운영비 차감 여부를 명시한다.}
\B{}{\small 프로그램 승인본(회사 정본 §5.5 B-01~B-08)의 연간 절감 214억 중 P1의 위치를 아래와 같이 구분한다. 달성 시점은 프로그램 착수 24개월 후(M+24, 2028-01)이며, 판정은 G5(2028-03-31)에서 한다.\par}
{\scriptsize
\begin{tabularx}{\linewidth}{|L{9mm}|L{22mm}|L{22mm}|L{20mm}|L{13mm}|L{34mm}|Y|Y|Y|}\hline
\hd{ID} & \hd{편익 항목} & \hd{현재 기준값} & \hd{목표값} & \hd{달성 시점} & \hd{연간 금액(억)} & \hd{산출 분모} & \hd{이 프로젝트 기여} & \hd{측정 책임자} \\ \hline
\ifwbblank
\rd{11mm}B-01 & \g{[편익 항목]} & \g{[값·단위]} & \g{[값·단위]} & \g{[M+n]} & \g{[금액 또는 ‘미산정(사유)’]} & \g{[분모와 출처]} & \g{[직접 / 간접 / 타 프로젝트 의존]} & \g{[직책·실명]} \\ \hline
\rd{11mm}B-02 & & & & & & & & \\ \hline
\rd{11mm}B-03 & & & & & & & & \\ \hline
\rd{11mm}B-04 & & & & & & & & \\ \hline
\rd{11mm}B-05 & & & & & & & & \\ \hline
\rd{11mm}B-06 & & & & & & & & \\ \hline
\rd{11mm}B-07 & & & & & & & & \\ \hline
\rd{11mm}B-08 & & & & & & & & \\ \hline
\rd{8mm}\textbf{소계} & \multicolumn{8}{l|}{\g{직접 기여 [\ \ \ \ ]억 / 의존 [\ \ \ \ ]억 / 미산정 [\ \ ]건}} \\ \hline
\rd{8mm}\textbf{차감} & \multicolumn{8}{l|}{\g{이관 후 연간 운영비 [\ \ \ \ ]억 (차감 여부 명시)}} \\ \hline
\else
B-04 & 폐기율 & 3.9\% & 2.2\% & M+24 & 43 & 식자재 취급원가 2,530억(C-15) × 1.7\%p & \textbf{직접} (통합 재고·수요 가시성) & 오정근 영업본부장 / 측정: 구매팀 폐기집계 \\ \hline
B-05 & 배달수수료 회피 & 배달 경유 242.7억 × 12.8\% & 앱 주문 전환 & M+24 & 37 & 소비자 대면 총매출 5,920억(C-14) 기반 & \textbf{직접} (주문 허브·앱 연동) — 단, 앱 개편은 자사앱 벤더 의존 & 오정근 영업본부장 / 측정: 배달 플랫폼 정산서 \\ \hline
B-07 & 매장 인건비율 & 21.4\% & 20.1\% & M+24 & 38 & 직영 소매 매출 2,980억(C-12) × 1.3\%p & \textbf{직접+P3 의존} (POS 운영 부하 감소는 P1, 정착은 P3 교육) & 오정근 영업본부장 / 측정: 부문별 손익 \\ \hline
B-06 & 물류 처리량 & 1,860 케이스/h & 4,200 케이스/h & M+24 & 96 & 이사회 자료에 분모 미기재 & \textbf{P2 의존} — P1은 재고 원장·인터페이스로 기여하나 처리량은 설비가 결정 & 한동석 물류본부장 / 측정: 센터운영 보고 \\ \hline
B-02 & 재고 정확도 & 91.3\% & 98.5\% & M+24 & (B-04에 기여, 별도 금액 없음) & 순환실사 SKU 69\% — \textbf{분모 확장 필요} & 직접 & 한동석 물류본부장 (분모 확장은 나영선 본부장 협의) \\ \hline
B-03 & 결품률 & 4.7\% & 1.5\% & M+24 & \textbf{미산정} — 기회손실을 현행 시스템으로 계산 불가. 신 플랫폼 가동 후 M+6에 산정식 확정 & 매장 기준 & 직접 & 오정근 영업본부장 \\ \hline
B-08 & 가맹 발주 리드타임 & 38시간 & 12시간 & M+12 & (B-03에 기여) & 가맹 발주 웹 로그 & 직접 (가맹 포털) & 오정근 영업본부장 / 측정: 포털 로그 \\ \hline
B-01 & 온라인 매출 비중 & 6.2\% & 18\% & M+24 & (B-05로 환산) & C-14 5,920억 & 직접+앱 벤더 의존 & 오정근 영업본부장 \\ \hline
\textbf{소계} & \multicolumn{8}{l|}{\textbf{직접 기여 118 / P2 의존 96 / 미산정 1건}} \\ \hline
\textbf{차감} & \multicolumn{8}{l|}{이관 후 연간 운영비 \textbf{−27} (P-27, 프로그램 전체 플랫폼 운영비. 승인본 214억에 미차감) — 달성 시점 M+17~, 책임자 박준영 서비스운영팀장} \\ \hline
\fi
\end{tabularx}}
\landscapeoff

% ------------------------------------------------------------------ 5
\secbar{5. 비용 (Costs) — 얼마가 드는가}
\guide{구축 비용(용역·라이선스·인프라·장비·전환·검증·PMO·예비비)과 운영 비용(이관 후 연간 × 기간)을 구분한다. 예산이 여러 층(승인 예산 / 집행 한도 / 계획 원가)이면 각 층의 정의를 밝히고, 5년 TCO를 한 줄로 요약한다. 단위 통일(억 원, 소수 첫째 자리).}
{\footnotesize
\begin{tabularx}{\linewidth}{|L{46mm}|L{22mm}|Y|}\hline
\hd{구축 비용 내역} & \hd{금액(억)} & \hd{산출 근거} \\ \hline
\ifwbblank
\rd{8mm}\g{[항목]} & & \g{[단가 × 수량 등]} \\ \hline
\rd{8mm} & & \\ \hline
\rd{8mm} & & \\ \hline
\rd{8mm} & & \\ \hline
\rd{8mm} & & \\ \hline
\rd{8mm} & & \\ \hline
\rd{8mm} & & \\ \hline
\rd{8mm}\g{예비비 (통제 주체 명시)} & & \\ \hline
\rd{8mm}\textbf{합계 (승인 예산)} & & \\ \hline
\rd{11mm}예산 계층 & \multicolumn{2}{l|}{\g{승인 예산 [\ \ \ \ ] / 집행 한도 [\ \ \ \ ] / 계획 원가 [\ \ \ \ ] — 각 층의 정의}} \\ \hline
\rd{11mm}연차별 집행 계획 & \multicolumn{2}{l|}{\g{[연도] [\ \ ] / [연도] [\ \ ] / [연도] [\ \ ] — 프로그램 연차 예산과 정합 확인}} \\ \hline
\rd{11mm}운영 비용 & \multicolumn{2}{l|}{\g{연 [\ \ \ ]억 × [n]년 = [\ \ \ \ ] (소관: [\ \ \ \ ])}} \\ \hline
\rd{11mm}\textbf{5년 TCO} & \multicolumn{2}{l|}{\g{구축 [\ \ \ \ ] + 운영 [\ \ \ \ ] = [\ \ \ \ ]}} \\ \hline
\else
SI 개발 용역 & 68.2 & FP 12,400점 × 55만원 (P-08b) \\ \hline
상용SW 라이선스 및 3년 유지보수 & 23.4 & API 게이트웨이 6.8 / MDM 9.2 / APM 4.1 / SAP 애드온 3.3 \\ \hline
인프라 & 23.9 & 클라우드 3년 선약정 14.6 / DR·백업 4.2 / 망분리·보안장비 5.1 \\ \hline
POS 단말·주변기기 & 8.6 & 직영 477대 × 148만원 = 7.06 + 설치 1.54 (가맹 단말 미포함) \\ \hline
데이터 전환 용역 & 11.4 & 1,842 테이블 / 7.4TB, SI와 별도 계약 \\ \hline
통합·성능시험 및 3자 검증 & 9.7 & 부하시험 3.1 / 보안진단 2.4 / 외부 감리 4.2 \\ \hline
발주사 PMO·형상관리·품질보증 & 12.6 & 상주 PMO 3명 × 24개월 + 도구 \\ \hline
P1 컨틴전시 (PM 통제) & 10.2 & 승인 예산의 6.07\% \\ \hline
\textbf{합계 (승인 예산)} & \textbf{168.0} & \\ \hline
예산 계층 & \multicolumn{2}{Y|}{승인 예산 168.0 / 집행 한도 156.0 (프로그램 관리예비비 12.0 유보, 스폰서 결재) / 계획 원가 145.8 (컨틴전시 10.2 제외) — 출처 P-08} \\ \hline
연차별 집행 계획 [추가 설정] & \multicolumn{2}{Y|}{2026년 102.0 / 2027년 60.0 / 2028년 6.0(하자보수·잔여) — 프로그램 연차 예산 176/198/54와 정합 확인 필요} \\ \hline
운영 비용 & \multicolumn{2}{Y|}{연 27.0억 × 3년 = 81.0 (P5 소관, 플랫폼 전체) — 출처 P-27} \\ \hline
\textbf{5년 TCO} & \multicolumn{2}{Y|}{구축 168.0 + 운영 27.0 × 4.5년(2027-06~2031-12) ≈ 121.5 = \textbf{약 289.5}} \\ \hline
\fi
\end{tabularx}}

% ------------------------------------------------------------------ 6
\secbar[\B{120mm}{60mm}]{6. 재무 평가 (Financial Case) — 계산이 맞는가}
\guide{편익과 비용을 같은 시간축에 놓고 회수기간·NPV·IRR 중 조직이 쓰는 지표를 적는다. 할인율과 편익 실현 기간을 명시하고, 편익이 일부만 실현될 때의 민감도를 한 줄 이상. 정확한 NPV보다 \textbf{“편익이 몇 \% 실현되어야 손익분기인가”}가 의사결정자에게 더 유용하다. 총편익에서 운영비를 뺀다.}
\B{}{\small [계산은 추가 설정, 승인본에는 NPV가 없음]\par}
{\footnotesize
\begin{tabularx}{\linewidth}{|L{40mm}|Y|Y|}\hline
\hd{지표} & \hd{값} & \hd{가정} \\ \hline
\ifwbblank
\rd{10mm}편익 실현 곡선 & \g{[연도별 실현률 \%]} & \g{[ramp-up 가정, 제외한 편익]} \\ \hline
\rd{10mm}운영비 & \g{[연차별]} & \g{[귀속 방식]} \\ \hline
\rd{10mm}연차별 순현금흐름(억) & & \\ \hline
\rd{10mm}회수기간 & \g{[착수 기준 n년 / 컷오버 기준 n년]} & \g{[편익 실현 시점·ramp-up 가정]} \\ \hline
\rd{10mm}NPV (할인율 [\ \ ]\%, [\ ]년) & & \g{[할인율 근거]} \\ \hline
\rd{10mm}민감도 & \g{편익 50\% / 70\% / 100\% 실현 시 [\ \ \ \ ]} & \\ \hline
\rd{10mm}\textbf{손익분기 편익 실현률} & \g{[\ \ ]\%} & \g{[NPV=0이 되는 실현률]} \\ \hline
\rd{10mm}연차별 현금 유출 & \g{[연도] / [연도] / [연도]} & \g{[연차 승인 예산과 대조]} \\ \hline
\else
편익 실현 곡선 & 2027년 30\%(컷오버 후 안정화) → 2028년 이후 100\% & P1 직접 기여 118억 기준. P2 의존 96억은 제외 \\ \hline
운영비 & 2027년 16억(6~12월), 2028년 이후 연 27억 & 플랫폼 운영비 전액을 P1에 귀속(보수적) \\ \hline
연차별 순현금흐름(억) & 2026 −102 / 2027 −41 / 2028 +85 / 2029 +91 / 2030 +91 / 2031 +91 & \\ \hline
회수기간 & 착수 기준 약 3.4년(2029년 중) / 컷오버 기준 약 2.2년 & 누적 −102 → −143 → −58 → +33 \\ \hline
NPV (할인율 8\%, 5년) & 약 +134억 & 할인율 = 가중평균 차입금리 5.2\%(C-06 산식) + 자기자본 프리미엄 [추가 설정] \\ \hline
민감도 & 편익 100\% 실현: NPV +134 / 70\%: 약 +16 / 50\%: 약 −63 (5년 내 미회수) & \\ \hline
\textbf{손익분기 편익 실현률} & \textbf{약 66\%} & 직접 기여 118억의 66\% ≈ 78억이 M+24부터 매년 실현되어야 5년 NPV ≥ 0 \\ \hline
연차별 현금 유출과 프로그램 예산 & 2026년 P1 102억은 프로그램 1차연도 176억의 58\% & CFO의 1차연도 148억 축소 요구 시 P1 배분 재검토 대상 \\ \hline
\fi
\end{tabularx}}
\begin{fieldbox}
\guide{요약 문단 한 개 — 의사결정자에게 전할 한 문장.}\par
\B{\lines{2}{9mm}}{편익의 3분의 1이 실현되지 않으면 이 프로젝트는 5년 안에 돈을 벌지 못한다. 이 숫자는 이사회에 보고된 적이 없으며, G1에서 스폰서에게 명시적으로 제시한다.}
\end{fieldbox}

% ------------------------------------------------------------------ 7
\secbar[\B{150mm}{60mm}]{7. 주요 리스크와 가정 (Key Risks \& Assumptions)}
\guide{편익·비용 숫자가 무너질 수 있는 상위 3~5개 리스크(원인 → 사건 → 편익/비용에의 영향)와, 그 숫자가 전제하는 가정(“~라고 가정한다. 이 가정이 깨지면 ~”). 리스크 등록부를 만들지 않는다 — 비즈니스 케이스가 틀릴 수 있는 이유만. 각 항목이 어느 숫자에 연결되는지 밝힌다.}
\begin{fieldbox}
\ifwbblank
\g{R1:} \lines{2}{8.5mm}
\g{R2:} \lines{2}{8.5mm}
\g{R3:} \lines{2}{8.5mm}
\g{R4:} \lines{2}{8.5mm}
\g{A1: [~라고 가정한다]. 이 가정이 깨지면 [\ \ ].} \lines{2}{8.5mm}
\g{A2:} \lines{2}{8.5mm}
\g{A3:} \lines{2}{8.5mm}
\else
\begin{itemize}
\item \textbf{R1:} 한동석 물류본부장이 처리량 4,200 케이스/h를 “현장에서 나온 숫자가 아니다”라고 공개 발언 → P2 편익 96억의 근거 재검토 → 프로그램 총편익 214억의 45\%가 흔들림. \textbf{P1 재무 평가는 이 96억을 이미 제외했다.}
\item \textbf{R2:} 자사앱 소스 형상관리를 외주 벤더가 보유 → 앱 개편이 P1 주문 허브 일정에 맞지 않음 → B-05 배달수수료 회피 37억과 B-01 온라인 비중이 M+24 이후로 지연.
\item \textbf{R3:} 가맹 계약 부속서 개정(388점 서면 동의, 2026-08-14 목표)이 지연 → 가맹점 POS SW 배포 범위 축소 → 가맹 판매 데이터 실시간 수집 불가 → B-02 재고 정확도·B-04 폐기율 목표 중 가맹 부분 미달.
\item \textbf{R4:} CFO의 1차연도 집행 148억 축소 요구가 확정 → P1 2026년 배분 축소 → 설계 스프린트 인력 축소 또는 라이선스 선약정 연기 → 컷오버 창(2027-02) 이탈 시 다음 창은 2027년 하반기.
\item \textbf{A1:} 편익 실현 곡선은 2027년 30\%, 2028년 이후 100\%를 가정한다. 이 가정이 깨지면(예: P3 교육 지연으로 활용률 저조) 회수기간이 1년 이상 늘어난다.
\item \textbf{A2:} 대성로지스 WMS 준실시간 연동은 물류본부의 3PL 계약 변경 협상이 2026년 상반기에 타결된다고 가정한다. 깨지면 통합 재고의 물류센터 부분은 야간 배치로 남고 B-02 목표 98.5\%는 매장 재고에 한정된다.
\item \textbf{A3:} 재고 정확도 목표 98.5\%의 분모는 현행 순환실사 SKU 69\%로 가정한다. 분모를 100\%로 확장하면 기준값 91.3\% 자체가 재측정 대상이다.
\end{itemize}
\fi
\end{fieldbox}

% ------------------------------------------------------------------ 8
\secbar[\B{125mm}{60mm}]{8. 권고와 승인 조건 (Recommendation \& Approval Conditions)}
\guide{채택 대안, 승인 요청 금액·기간, 승인에 붙는 조건(번호를 붙여 헌장·게이트 문서에서 “승인 조건 3항”으로 인용 가능하게), 다음 검토 시점(게이트). 이사회·투자심의가 붙인 조건은 그대로 옮긴다.}
\begin{fieldbox}
\ifwbblank
\g{권고: [채택 대안, 승인 요청 금액·기간]} \lines{2}{8.5mm}
\g{승인 조건 (원문 전재)} \par
\g{1.} \lines{1}{8.5mm} \g{2.} \lines{1}{8.5mm} \g{3.} \lines{1}{8.5mm} \g{4.} \lines{1}{8.5mm} \g{5.} \lines{1}{8.5mm}
\g{PM 추가 조건 (있으면)} \lines{1}{8.5mm}
\g{다음 검토: [게이트, 날짜, 검토 항목]} \lines{1}{8.5mm}
\else
\begin{itemize}
\item \textbf{권고:} 대안 3(통합 플랫폼 SI 구축)을 승인 예산 168.0억, 기간 2026-01-05 ~ 2027-05-29로 추진한다. 대안 1(최소 조치)은 2026-06까지 폴백 옵션으로 유지하고 G1에서 폐기 여부를 결정한다.
\item \textbf{이사회 승인 조건 (2025-12 의결문 전재):}
  \begin{enumerate}[leftmargin=6mm]
  \item 총예산 428억은 3개년(176 / 198 / 54)에 걸쳐 집행하며, 연차별 집행은 재무본부 자금 계획과 연동해 재확인한다.
  \item 프로그램 관리예비비 24억은 원천 유보하며 집행은 프로그램 스폰서 결재 사항으로 한다.
  \item 6개 게이트(G0~G5)로 단계별 승인하며, 게이트마다 편익 실현 전망을 재확인한다.
  \item M+24 연간 절감 214억의 실현 여부를 G5(2028-03-31)에서 판정한다.
  \item 2020년·2023년 내부감사 지적(편익 산정 근거 부재, 손실 확정 절차 부재, 위탁 산출물 형상관리·소스 인수 절차 미비)을 반복하지 않는 관리 체계를 갖춘다.
  \end{enumerate}
\item \textbf{P1 추가 조건 (PM 제안, G1 확정 대상):} 6. P1 편익은 직접 기여 118억으로 한정해 관리하고 B-06 96억은 P2 편익으로 이관한다. 7. 결품률 기회손실 산정식을 M+6까지 확정한다.
\item \textbf{다음 검토:} G1 2026-04-30 — 원가·일정 기준선, FP 기준선 확인서, 편익 실현률 손익분기(66\%) 재산정.
\end{itemize}
\fi
\end{fieldbox}

% ------------------------------------------------------------------ 9
\secbar[\B{110mm}{60mm}]{9. 편익 책임자와 검토 일정 (Benefit Owners \& Review Schedule)}
\guide{편익 항목별 실현 책임자(현업 임원 — 부서명이 아니라 \textbf{직책과 실명}), 측정 담당(데이터를 뽑는 사람), 데이터 출처, 검토 시점(프로젝트 종료 후의 날짜). PM을 편익 책임자로 적지 않는다 — PM은 종료 후 그 자리에 없다.}
{\footnotesize
\begin{tabularx}{\linewidth}{|L{34mm}|L{34mm}|L{30mm}|Y|L{30mm}|}\hline
\hd{편익 ID} & \hd{실현 책임자(직책·실명)} & \hd{측정 담당(직책·실명)} & \hd{데이터 출처} & \hd{검토 시점} \\ \hline
\ifwbblank
\rd{9mm}\g{[B-0n 항목]} & & & & \g{[M+n, G5 등]} \\ \hline
\rd{9mm} & & & & \\ \hline
\rd{9mm} & & & & \\ \hline
\rd{9mm} & & & & \\ \hline
\rd{9mm} & & & & \\ \hline
\else
B-04 폐기율 & 오정근 영업본부장 & 강현도 구매팀장 & 구매팀 폐기집계 (월) & M+12 (2027-01), G5 (2028-03-31) \\ \hline
B-05 배달수수료 회피 / B-01 온라인 비중 & 오정근 영업본부장 & 재경팀장 & 배달 플랫폼 정산서, 부문별 손익 & M+12, G5 \\ \hline
B-07 매장 인건비율 & 오정근 영업본부장 & 인사본부 급여 담당 & 부문별 손익 (분모 C-12) & M+12, G5 \\ \hline
B-02 재고 정확도 & 한동석 물류본부장 & 물류본부 순환실사 담당 & 순환실사 보고 — 분모 확장 방안은 나영선 본부장과 M+3까지 합의 & M+12, G5 \\ \hline
B-03 결품률 / B-08 발주 리드타임 & 오정근 영업본부장 & 영업본부 매장운영 KPI 담당 & 매장운영 KPI, 가맹 포털 로그 & M+6 산정식 확정, M+12, G5 \\ \hline
B-06 물류 처리량 & 한동석 물류본부장 (P2 소관) & 물류본부 센터운영 담당 & 센터운영 보고 & P2 편익 계획으로 이관 \\ \hline
\fi
\end{tabularx}}

\needspace{26mm}\vspace{3mm}
{\footnotesize
\begin{tabularx}{\linewidth}{|L{30mm}|Y|L{28mm}|L{40mm}|}\hline
\hd{확인} & \hd{직책·실명} & \hd{날짜} & \hd{서명} \\ \hline
\rd{8mm}작성 (PM) & \B{}{P1 PM (발주사 IT본부)} & \B{}{2026-01-05} & \\ \hline
\rd{8mm}검토 (스폰서 / 재무) & \B{}{정미란 CFO (프로그램 스폰서) / 재경팀장} & \B{}{2026-01-05} & \\ \hline
\end{tabularx}}

\end{document}
